\newpage
\begin{center}
  \textbf{\large ЗАКЛЮЧЕНИЕ}
\end{center}
\refstepcounter{chapter}
\addcontentsline{toc}{chapter}{ЗАКЛЮЧЕНИЕ}

В результате выполнения выпускной квалификационной работы был проведен анализ 
существующих методов визуальной и лидарной одометрии, а также мультимодальных
подходов. Выделены основные классы широко используемых методов. В качестве основы 
для разработки был выбран полностью нейросетевой метод DVLO. 

Идея метода DVLO состоит в особой процедуре экстракции, а также объединения локальных и
глобальных признаков облаков точек и изображений. В частности, используется подход кластеризации
признаков, предложенный в работе Image as Set of Points.

Используя открытый набор данных KITTI Odometry были проведены эксперименты с базовой версией метода,
а также с предложенной модификацией, использующей вторую камеру для уточнения предсказаний.
После обучения и валидации удалось показать значительное падение средних значений ошибок
положения и ориентации (на $30\%$ и $26\%$ соответственно).

Полученные характеристики работы разработанного алгоритма соответствуют требованиям, 
поставленным в техническом задании --- частота работы не менее 8 Гц, количество доступной 
видеопамяти не менее 8 Гб, относительные ошибки расстояния и ориентации менее 2\% и 0.6$^\circ/100$м.

Для проведения сравнительного анализа также были поставлены эксперименты с классическими 
методами визуальной и лидарной одометрии. Результаты свидетельствуют о превосходстве
современных классических методов, основанных на оптимизации над полностью нейросетевыми.
Однако, последние обладают потенциалом к существенному расширению области применения
за счет возможности дообучения параллельных слоев для решения других задач. Кроме того,
нейросетевые модели могут быть боле робастными при обучении на разнообразных данных, в то время
как классические методы часто требуют тщательной настройки параметров перед использованием в конкретном окружении.

В данной работе продемонстрирована возможность масштабирования нейросетевых методов на новые сенсоры, что
указывает на перспективность подобных разработок в будущем.

